\subsection*{3.1 Random hypergraphs mix well}
In this section, we'll show that given an expanding graph $G$, its random hypergraph completion is, with high probability, a good expander. Throughout, we say that an undirected regular graph $G$ is a $\lambda$-expander if the second largest eigenvalue of its normalized adjacency matrix, in magnitude, is at most $\lambda$.
\paragraph{Definition17 (random hypergraph completion of a graph).} Let $G=(V=[n],E_G)$ be a $d$-regular graph. To sample a random hypergraph completion, $\mathbf H$ of $G$, we choose a uniformly random ordering of $G$'s edges, $\{(u_1,v_1),\ldots,(u_{nd},v_{nd})\}=E_G$. Then, we set $\mathbf H=(V,E_{\mathbf H})$, where
\[E_{\mathbf H}\coloneqq\{(u_i,\lceil i/d\rceil,v_i)\mid i\in[nd]\}.\]
Equivalently, for each $(u,v)\in E_G$, we choose an independent and uniform $\mathbf w\in V$ and add the hyperedge $(u,\mathbf w,v)$ to $E_{\mathbf H}$, conditioned on the resulting hypergraph $\mathbf H$ being $d$-regular.
Next, we prove that the random hypergraph completion mixes well with high probability.
\paragraph{Lemma18.} Let $G=(V=[n],E)$ be a $d$-regular $\lambda$-expander and $\mathbf H$ be a random hypergraph completion of $G$. With probability at least $1-2^{-n}$, $\mathbf H$ is a $\lambda'$-mixing hypergraph for $\lambda'=2\lambda+2/\sqrt d$.
In order to prove Lemma18, we'll use Hoeffding's inequality for sampling without replacements.
\paragraph{Fact19 (Hoeffding's inequality,[18]).} For any integers $a,k\leq m$, suppose there are $m$ items, of which $k$ of them are marked. Let $\mathbf x$ be a random variable indicating the number of marked items when $a$ of the $m$ items are sampled uniformly and independently without replacement. Then, for any $t\geq0$,
\[\Pr[|\mathbf x-(k/m)\cdot a|\geq t]\leq2\exp(-2t^2/a).\]
\begin{proof}[Proof of Lemma18]
Fix arbitrary $S_1,S_2,S_3\subseteq V$. We will show that with probability at least $1-2^{-4n}$ it holds that
\[\left|\frac{E_{\mathbf H}(S_1,S_2,S_3)}d-\frac{|S_1|\cdot|S_2|\cdot|S_3|}{n^2}\right|\leq\left(\frac2{\sqrt d}+2\lambda\right)\sqrt{|S_1|\cdot|S_3|}\tag{5}\]
where $E_{\mathbf H}(S_1,S_2,S_3)$ is the number of edges $(v_1,v_2,v_3)$ in $\mathbf H$ where $v_j\in S_j$ for each $j\in[3]$. The desired result then follows from a union bound over the $(2^n)^3=2^{3n}$ choices for $S_1,S_2,S_3$.
Let $E_G(S_1,S_3)$ be the number of edges, $(u,v)$, of $G$, such that $u\in S_1$ and $v\in S_3$. By the expander mixing lemma applied to $G$, we have that
\[\left|\frac{E_G(S_1,S_3)}d-\frac{|S_1||S_3|}n\right|\leq\lambda\sqrt{|S_1||S_3|}.\]
Let us define $\mu\coloneqq|S_1||S_3|/n$ and $\Delta\coloneqq\lambda\sqrt{|S_1||S_3|}$. We consider two cases.
\begin{enumerate}
\item In the first case, $\mu\leq\Delta$. Here, we use the simple bound
\[0\leq\frac{E_{\mathbf H}(S_1,S_2,S_3)}d\leq\frac{E_G(S_1,S_3)}d\leq2\Delta\]
that holds with probability1. This, along with the fact $|S_1|\cdot|S_2|\cdot|S_3|/n^2\leq\mu\leq\Delta$ implies that Equation(5) always holds.
\item In the second case, $\mu>\Delta$. Here, we will apply Fact19. $G$ has a total of $nd$ edges, of which $d\cdot|S_2|$ are matched to a vertex in $S_2$. $E_{\mathbf H}(S_1,S_2,S_3)$ samples $E_G(S_1,S_3)$ of those $nd$ edges (without replacement) and counts how many were among the $d\cdot|S_2|$ assigned to $S_2$. Hence, by Hoeffding's inequality, we have for any $t\geq0$,
\[\Pr\left[\left|E_{\mathbf H}(S_1,S_2,S_3)-\frac{d\cdot|S_2|}{nd}\cdot E_G(S_1,S_3)\right|\geq t\right]\leq2\exp\left(\frac{-2t^2}{E_G(S_1,S_3)}\right).\]
Then, setting $t=2\sqrt{d|S_1|\cdot|S_3|}$ and using the fact that $E_G(S_1,S_3)\leq d(\mu+\Delta)\leq2d\mu$,
\begin{multline*}
\Pr\left[\left|E_{\mathbf H}(S_1,S_2,S_3)-\frac{|S_2|}n\cdot E_G(S_1,S_3)\right|\geq2\sqrt{d|S_1|\cdot|S_3|}\right]\\
\leq2\exp\left(\frac{-8d|S_1|\cdot|S_3|}{2d\cdot(|S_1||S_3|/n)}\right)=2\exp(-4.5n)\leq2^{-4n}.
\end{multline*}
As the above shows, $E_{\mathbf H}(S_1,S_2,S_3)$ is within $\pm t$ of its expectation with probability at least $2^{-4n}$. When that occurs, we have that
\begin{align*}
&\left|\frac{E_{\mathbf H}(S_1,S_2,S_3)}d-\frac{|S_1|\cdot|S_2|\cdot|S_3|}{n^2}\right|\\
&=\frac1d\left|E_{\mathbf H}(S_1,S_2,S_3)-\frac{|S_2|d\mu}n\right|\\
&\leq\frac td+\frac{|S_2|}n\left|\frac{E_G(S_1,S_3)}d-\mu\right|\\
&\leq\frac td+\frac{|S_2|}n\cdot\Delta\qquad\text{(expander mixing lemma)}\\
&\leq2\sqrt{\frac{|S_1|\cdot|S_3|}d}+\lambda\sqrt{|S_1|\cdot|S_3|}.\qquad(|S_2|/n\leq1)
\end{align*}
\end{enumerate}
Hence, Equation(5) holds with probability at least $1-2^{-4n}$ in both cases, so we can union bound over the $2^{3n}$ choices for $S_1,S_2,S_3$.
\end{proof}
For our purposes, we will want the hypergraph to be symmetric. It is quite easy to ``symmetrize'' any hypergraph at only a modest cost to the degree.
\paragraph{Proposition20.} For any $n,d\in\mathbb N$, there is an algorithm running in time $O(nd)$ that takes as input any $d$-regular $\lambda$-mixing hypergraph $H=(V=[n],E_H)$ and outputs a $2d$-regular $\lambda$-mixing hypergraph $H'=(V=[n],E_{H'})$ over the same vertices.
\begin{proof}
For every edge $(u,v,w)\in E_H$ we include both $(u,v,w)$ and the reverse edge $(w,v,u)$ in $E_{H'}$.\footnote{Note that, we do this even if it results in duplicated hyperedges: If $(u,v,w)$ and $(w,v,u)$ are both in $E_H$, then there will be two copies of $(u,v,w)$ and two copies of $(w,v,u)$ in $E_{H'}$.} Clearly, this results in the degree of $H'$ being $2d$. Then, for any $S_1,S_2,S_3\subseteq V$,
\begin{align*}
&\left|\frac{E_{H'}(S_1,S_2,S_3)}{2d}-\frac{|S_1|\cdot|S_2|\cdot|S_3|}{n^2}\right|\\
&=\left|\frac{E_H(S_1,S_2,S_3)+E_H(S_3,S_2,S_1)}{2d}-\frac{|S_1|\cdot|S_2|\cdot|S_3|}{n^2}\right|\\
&\leq\frac12\left|\frac{E_H(S_1,S_2,S_3)}d-\frac{|S_1|\cdot|S_2|\cdot|S_3|}{n^2}\right|\\
&\quad+\frac12\left|\frac{E_H(S_3,S_2,S_1)}d-\frac{|S_1|\cdot|S_2|\cdot|S_3|}{n^2}\right|\\
&\leq\lambda\sqrt{|S_1|\cdot|S_3|}.
\end{align*}
Therefore, $H'$ is $\lambda$-mixing, as desired.
\end{proof}
We have explicit (and even fully explicit) constructions of Ramanujan graphs, i.e., $\lambda$-expanders for $\lambda\leq2\sqrt{d-1}/d$, albeit with some restrictions on $d$[26,28]. By manipulating Ramanujan graphs, Alon gave a construction of $d$-regular $\lambda$-expanders over $n$ vertices for any $d$ and $n$ while suffering only a tiny loss in $\lambda$ (see[1], and also[29,13] for weaker constructions). In particular, there exist explicit expanders with $\lambda=O(1/\sqrt d)$ for all $n$-s. We thus get the following corollary.
\paragraph{Corollary21.} There exists a probabilistic algorithm such that for any integer $n$ and even integer $6\leq d\leq n$, runs in time $\operatorname{poly}(n)$ and with probability at least $1-2^{-n}$ outputs a3-uniform symmetric $d$-regular hypergraph that is $\lambda=c_{\mathrm{rand}}/\sqrt d$-mixing, where $c_{\mathrm{rand}}\geq2$ is some universal constant.
We will refer to the above probabilistic construction as our preprocessing step.
